\documentclass[a4paper,11pt]{article}

\usepackage{amsmath, amssymb}
\usepackage{geometry}
\geometry{margin=18mm}

\usepackage{hyperref}
\usepackage{setspace}
\setstretch{0.95}

\title{Create or Curate?\\ How Instructor Framing Shapes the Value of Educational Video\\
\large  1022179 Riki Nishino\\
Advisor: Ian Frank}
\author{}
\date{}

\begin{document}

\maketitle

\section{Research Background}

In recent years, the use of educational videos has expanded in higher education due to the promotion of digital transformation (DX) and the widespread adoption of online learning environments. Noetel et al. (2021) conducted a systematic review and meta-analysis and found that educational videos improve learning outcomes in higher education. These findings indicate that educational videos have become valuable instructional resources in university education.
However, creating educational videos requires considerable time, technical skills, and digital competence. Liesa-Orús et al. (2023) reported substantial differences in the digital competence of university instructors, suggesting that not all instructors are able to produce their own educational videos. Consequently, selecting and utilizing existing educational videos has become an important instructional practice alongside creating new materials.
In this context, it is important to consider not only educational videos themselves but also how instructors present them to students. Even when the same video is used, students may perceive it differently depending on whether it is described as an instructor-created resource or as a carefully curated resource. Despite this possibility, the influence of such instructor-provided information on students' evaluations of educational videos has received little empirical attention.

\section{Research Purpose}

The purpose of this study is to examine how instructor-provided information attached to educational videos influences learners' evaluations of those videos.
Specifically, this study investigates whether learners' perceptions of educational value, usefulness, trustworthiness, satisfaction, and learning motivation differ when the same  is described either as having been created by the instructor or as having been carefully curated by the instructor.
By doing so, this study aims to clarify the effectiveness of the concept of curation in education.

\section{Research Question}

How effective is the concept of curation in education?
More specifically, how do different descriptions of the same educational video influence learners' evaluations of the video?

\section{Related Work}

Previous studies have shown that instructors commonly both create and select educational videos for instructional purposes. Shoufan and Mohamed (2022) reviewed the use of YouTube in education and argued that instructors not only produce their own videos but also integrate existing videos into their teaching. In addition, Deschaine and Sharma (2015) defined digital curation as the educational practice of selecting, organizing, and sharing learning resources to support teaching and learning. Flintoff et al. (2014) also emphasized the importance of digital curation in supporting education, research, and professional development.
Research in communication has long demonstrated that people's evaluations are influenced not only by information itself but also by their perceptions of its source. Hovland and Weiss (1951) showed that source credibility significantly affects how identical information is received and evaluated.
More recent studies in online education suggest that learners' evaluations can also be influenced by instructor-related information and presentation cues, even when the instructional content remains the same. Ramlatchan and Watson (2020) found that different presentation formats of online educational videos affected learners' perceptions of instructor credibility and psychological closeness. Similarly, Bateman and Schmidt-Borcherding (2018) argued that various signals in educational videos, including visual and design elements, shape learners' interpretation and evaluation.
However, few studies have examined whether merely describing the same educational video as either instructor-created or instructor-curated influences learners' evaluations of that video.

\section{Research Method}

This study will conduct a small-scale experiment using the same educational video.
Participants will be randomly assigned to one of three conditions:

\begin{itemize}
  \item \textbf{Creation condition:}“This video was created by the instructor for this course.”
  \item \textbf{Curation condition:}“This video was curated by the instructor for this course.”
  \item \textbf{Neutral condition:}No additional information about the origin of the video is provided.
\end{itemize}

To isolate the effect of instructor-provided information, the same video will be used in all conditions.
After viewing the video, participants will complete a comprehension test and a questionnaire. The questionnaire will measure perceived educational value, usefulness, trustworthiness, satisfaction, and learning motivation.
The results will be compared across conditions to analyze how instructor framing influences learners’ perceptions and evaluations of educational videos.

\section{Expected Contributions}

This study does not compare the quality or learning effectiveness of different educational videos themselves, but rather focuses on how instructor framing influences learners’ perceptions.
The findings are expected to contribute to understanding the value of selecting and appropriately contextualizing existing videos in situations where producing new educational videos is difficult.
Furthermore, this study aims to provide evidence that in the use of educational videos, not only creating educational videos but also carefully selecting and contextualizing existing ones has educational value.


\begin{thebibliography}{99}

\bibitem{noetel2021}
Noetel, M., et al. (2021).
Video improves learning in higher education: A systematic review.
\textit{Review of Educational Research}, 91(2), 204--236.

\bibitem{liesa2023}
Liesa-Orús, M., et al. (2023).
Digital competence of university teachers: A systematic review.
\textit{Education Sciences}, 13(5), 508.

\bibitem{shoufan2022}
Shoufan, A., \& Mohamed, F. (2022).
YouTube and Education: A Scoping Review.
\textit{IEEE Access}, 10, 125547--125592.

\bibitem{deschaine2015}
Deschaine, M. E., \& Sharma, R. (2015).
The five stages of digital curation: Supporting teaching and learning in online learning environments.
\textit{Journal of Online Learning and Teaching}, 11(3), 398--409.

\bibitem{flintoff2014}
Flintoff, K., Mellow, P., \& Clark, M. (2014).
Digital curation: Opportunities for learning, teaching, research and professional development.
In \textit{Proceedings of ASCILITE 2014}.

\bibitem{hovland1951}
Hovland, C. I., \& Weiss, W. (1951).
The influence of source credibility on communication effectiveness.
\textit{Public Opinion Quarterly}, 15(4), 635--650.

\bibitem{ramlatchan2020}
Ramlatchan, M., \& Watson, G. S. (2020).
Enhancing instructor credibility and immediacy in online multimedia designs.
\textit{Educational Technology Research and Development}, 68, 3239--3264.

\bibitem{bateman2018}
Bateman, J. A., \& Schmidt-Borcherding, F. (2018).
The communicative effectiveness of education videos:
Towards an empirically motivated multimodal account.
\textit{Multimodal Technologies and Interaction}, 2(3), 59.

\end{thebibliography}

\end{document}